% GENERATED FILE. Edit canonical lessons, then run npm run generate:books.
% canonical-source-hash: fnv1a64:b47128d60123e9ed
% canonical-lessons: IT-C33-questo, IT-C33-quello, IT-C33-questo-caffe, IT-C33-questo-e, IT-C33-qui-li

\chapter{This and That}
\label{ch:it-questo-quello}
\hlchaptermodality{33}

\begin{chapteropening}
\textbf{By the end of this chapter:} \emph{I can point at a thing near me and a thing away from me, put a demonstrative in front of a noun, introduce a third person, and say here and there.}

The chapter's last lesson closes the pointing set and shows that ille — the Latin far-pointer — has now done three jobs in this book: the article of Chapter 1, the quello of two lessons back, and lì.
\end{chapteropening}

\section[questo]{questo — this one, here}

\label{lesson:IT-C33-questo}



\begin{quote}
A new chapter, opening a column this book has left entirely empty.

\begin{itemize}
\raggedright
  \item \emph{Recall:} say \emph{per + infinito} — \textbf{R1}, one lesson back, before it decays
  \item \emph{Recall:} say \emph{ma} — \textbf{R2}, the first real retrieval, five to fifteen lessons back
  \item \emph{Recall:} say \emph{prima il primo, dopo il secondo} — \textbf{R3}, durable, twenty to sixty lessons back
  \item \emph{Recall:} say \emph{sì} — \textbf{R3}, durable again, a second word from that distance
  \item \emph{Recall:} say \emph{il pane} — \textbf{R4}, recognition at distance, eighty lessons back or more
\end{itemize}
\end{quote}

\subsection*{You'll want to know: The word}
\begin{quote}
\textbf{questo} (masculine) · \textbf{questa} (feminine) — \textbf{this}. plural: \textbf{questi} · \textbf{queste} — \textbf{these}.
\end{quote}

\begin{quote}
\textbf{Questo è buono.} — This is good. \textbf{Questa è la mia amica.} — This is my friend. \emph{(the possessive is the next chapter's)} \textbf{Prendo questo.} — I'll have this one.
\end{quote}

Nothing in the previous thirty-two chapters could point at anything. You could name a coffee, and you could not indicate \emph{that} coffee rather than another.

It changes its ending like any describing word, and it can either stand alone — \emph{prendo \textbf{questo}} — or go in front of a noun, which is the next lesson.

\begin{cousinweb}
Latin's pointing words were \textbf{hic} (this, near me), \textbf{iste} (that, near you) and \textbf{ille} (that, over there). Spoken Latin found them too small to hear and reinforced them with \textbf{eccum}, \textquotedblleft{}behold, look\textquotedblright{} — the same gesture-word that is inside \emph{ecce homo}.

\textbf{Eccu(m) istum} $\to$ \textbf{questo}. Literally: \emph{look, that one of yours.}

You have met that reinforcing \emph{eccu-} twice already without being told. \textbf{Così} is \emph{eccu(m) sīc}, \textquotedblleft{}behold, thus.\textquotedblright{} And the article \textbf{il, la, lo} is what is left of \textbf{ille} — the third of those three Latin pointers, worn down until it stopped pointing at all and became \textquotedblleft{}the\textquotedblright{}. The next lesson takes \emph{ille} back the other way.
\end{cousinweb}

\subsection*{Your turn}
Say these aloud:

\begin{itemize}
\raggedright
  \item \emph{questo, questa}
  \item \emph{Questo è buono.}
  \item \emph{Prendo questo.}
\end{itemize}

\begin{tcolorbox}[breakable,colback=teal!4,colframe=teal!35!black,title={Before you move on}]
How do you say \textquotedblleft{}this\textquotedblright{}? (\emph{\textbf{Questo}}, or \emph{questa}.) Which Latin words is it made of? (\emph{\textbf{Eccu(m) istum}} — \textquotedblleft{}behold, that of yours\textquotedblright{}.) Which word already in this book carries the same \emph{eccu-}? (\emph{\textbf{Così}}.)
\end{tcolorbox}
\section[quello]{quello — that one, over there}

\label{lesson:IT-C33-quello}



\begin{quote}
One lesson back, and this chapter is already half built.

\begin{itemize}
\raggedright
  \item \emph{Recall:} say \emph{questo} — \textbf{R1}, one lesson back, before it decays
  \item \emph{Recall:} say \emph{o} — \textbf{R2}, the first real retrieval, five to fifteen lessons back
  \item \emph{Recall:} say \emph{no} — \textbf{R3}, durable, twenty to sixty lessons back
  \item \emph{Recall:} say \emph{l'acqua, il vino} — \textbf{R4}, recognition at distance, eighty lessons back or more
\end{itemize}
\end{quote}

\subsection*{You'll want to know: the pair, and it is only a pair}
\begin{quote}
\textbf{quello} (masculine) · \textbf{quella} (feminine) — \textbf{that}. plural: \textbf{quelli} · \textbf{quelle} — \textbf{those}.
\end{quote}

\begin{quote}
\textbf{Questo è buono, quello non è buono.} — This one's good, that one isn't. \textbf{Prendo quello.} — I'll have that one.
\end{quote}

\textbf{Italian points two ways, not three.} \emph{Questo} is near the speaker; \emph{quello} is everything else — near you, or far from both of us. Spanish keeps all three of Latin's — \emph{este}, \emph{ese}, \emph{aquel} — and a learner arriving from Spanish has one distinction to forget rather than one to learn.

That is the whole system. There is no third word.

\begin{cousinweb}
\textbf{Quello} is \textbf{eccu(m) illum} — \textquotedblleft{}behold, that one over there.\textquotedblright{}

\textbf{Illum} is \textbf{ille}, and you have owned \emph{ille} since the first chapter of this book. It is what \textbf{il}, \textbf{la} and \textbf{lo} are made of: the Latin far-pointer, used so often to mean \emph{the one already mentioned} that it stopped pointing and became the definite article.

So Italian took one Latin word down two roads. Unstressed and worn to a stump, it became \textbf{il} — a word that points at nothing. Reinforced with \emph{eccum} and kept strong, it became \textbf{quello} — a word that points hard.

Before a noun, \emph{quello} takes exactly the shapes the article takes, for exactly the same reason: \textbf{quel giorno} (\emph{il} giorno), \textbf{quello zucchero} (\emph{lo} zucchero), \textbf{quell'amico} (\emph{l'}amico), \textbf{quella sera} (\emph{la} sera). If you know the article rule you already know this one.
\end{cousinweb}

\subsection*{Your turn}
Say these aloud:

\begin{itemize}
\raggedright
  \item \emph{quello, quella}
  \item \emph{Questo è buono, quello non è buono.}
  \item \emph{quel giorno, quello zucchero, quell'amico}
\end{itemize}

\begin{tcolorbox}[breakable,colback=teal!4,colframe=teal!35!black,title={Before you move on}]
How many ways does Italian point? (\textbf{Two} — \emph{questo} and \emph{quello}.) Which Latin word is \emph{quello} built on? (\emph{\textbf{Ille}}.) And which word from the first chapter is built on the same \emph{ille}? (\textbf{The article} \emph{\textbf{il, la, lo}}.)
\end{tcolorbox}
\section[questo caffè]{questo caffè — where the pointing word goes, and what it evicts}

\label{lesson:IT-C33-questo-caffe}



\begin{quote}
The chapter's middle lesson, and it introduces no new word.

\begin{itemize}
\raggedright
  \item \emph{Recall:} say \emph{quello} — \textbf{R1}, one lesson back, before it decays
  \item \emph{Recall:} say \emph{perché} — \textbf{R2}, the first real retrieval, five to fifteen lessons back
  \item \emph{Recall:} say \emph{sì, grazie / no, grazie} — \textbf{R3}, durable, twenty to sixty lessons back
  \item \emph{Recall:} say \emph{undici — sedici} — \textbf{R4}, recognition at distance, eighty lessons back or more
\end{itemize}
\end{quote}

\subsection*{You'll want to know: in front, and instead of the article}
\begin{quote}
\textbf{questo caffè} — this coffee · \textbf{quella persona} — that person
\end{quote}

Two rules, and the second is the one to watch.

\textbf{One: it goes in front.} Italian's ordinary describing words go \emph{after} the noun — \emph{un vino rosso}, a red wine — but \emph{questo} and \emph{quello} go before it, always, exactly as English does.

\textbf{Two: it replaces the article.} From its first chapter this book has told you to store every noun with its article — \emph{il caffè}, \emph{la persona}, \emph{lo zucchero}. When a demonstrative arrives, the article goes:

\begin{quote}
\textbf{Right:} questo caffè · quella persona \textbf{Wrong:} \emph{questo il caffè} · \emph{il questo caffè} · \emph{quella la persona}
\end{quote}

Italian will not carry two of these at once. English does the same without comment — you cannot say \emph{the this coffee} — so the rule costs an English speaker nothing to obey and is worth stating because it is the only place in this book where a taught article has to be taken back out.

\textbf{This lesson introduces no new word.} \emph{Questo} and \emph{quello} are two and one lessons back; \emph{caffè}, \emph{persona} and \emph{zucchero} are Chapters 22 and 23.

\subsection*{Your turn}
Say these aloud:

\begin{itemize}
\raggedright
  \item \emph{questo caffè}
  \item \emph{quella persona}
  \item \emph{quello zucchero}
  \item \emph{questa amica}
\end{itemize}

\begin{tcolorbox}[breakable,colback=teal!4,colframe=teal!35!black,title={Before you move on}]
Where does \emph{questo} go? (\textbf{In front of the noun}.) What happens to the article? (\textbf{It goes} — \emph{questo caffè}, never \emph{questo il caffè}.) Which describing words go AFTER the noun instead? (\textbf{The ordinary ones} — \emph{un vino rosso}.)
\end{tcolorbox}
\section[Questo è Marco]{Questo è Marco — introducing somebody}

\label{lesson:IT-C33-questo-e}



\begin{quote}
Four in, and this one puts three chapters together.

\begin{itemize}
\raggedright
  \item \emph{Recall:} say \emph{questo caffè} — \textbf{R1}, one lesson back, before it decays
  \item \emph{Recall:} say \emph{quando} — \textbf{R2}, the first real retrieval, five to fifteen lessons back
  \item \emph{Recall:} say \emph{d'accordo} — \textbf{R3}, durable, twenty to sixty lessons back
  \item \emph{Recall:} say \emph{diciassette — venti} — \textbf{R4}, recognition at distance, eighty lessons back or more
\end{itemize}
\end{quote}

\subsection*{You'll want to know: the introduction}
\begin{quote}
\textbf{Questo è Marco.} — This is Marco. \textbf{Questa è Giulia.} — This is Giulia.
\end{quote}

\textbf{Questo} for a man, \textbf{questa} for a woman. The word agrees with the person being introduced, not with you and not with the person you are speaking to.

The \emph{mi chiamo} chapter taught you the whole of introducing \textbf{yourself} — \emph{mi chiamo}, \emph{come ti chiami}, \emph{piacere} — and left the third person out entirely, because a third-person introduction needs a pointing word and there was none. There is now.

The full exchange, using only what this book has given you:

\begin{quote}
— \textbf{Questo è Marco.} — \textbf{Piacere.} — \textbf{Piacere.}
\end{quote}

You can add what the person is to you, once the possessives arrive. For now, the plain noun does it:

\begin{quote}
\textbf{Questo è l'amico di Giulia.} — This is Giulia's friend.
\end{quote}

And note what Italian does not do. English says \emph{this is my friend Marco} and can also say \emph{meet Marco}. Italian has no imperative in this book yet, so \emph{questo è} is the introduction, and it is the ordinary one — an Italian would say it too.

\subsection*{Your turn}
Say these aloud:

\begin{itemize}
\raggedright
  \item \emph{Questo è Marco.}
  \item \emph{Questa è Giulia.}
  \item \emph{Questo è l'amico di Giulia.}
\end{itemize}

\begin{tcolorbox}[breakable,colback=teal!4,colframe=teal!35!black,title={Before you move on}]
How do you introduce a man? (\emph{\textbf{Questo è …}}) And a woman? (\emph{\textbf{Questa è …}}) Which chapter taught you to introduce yourself and could not teach this? (\textbf{The} \emph{\textbf{mi chiamo}} \textbf{chapter} — it had no word for \emph{this}.)
\end{tcolorbox}
\section[qui / lì]{qui, lì — here and there}

\label{lesson:IT-C33-qui-li}



\begin{quote}
The chapter closes where it began, with pointing.

\begin{itemize}
\raggedright
  \item \emph{Recall:} say \emph{Questo è Marco} — \textbf{R1}, one lesson back, before it decays
  \item \emph{Recall:} say \emph{per + infinito} — \textbf{R2}, the first real retrieval, five to fifteen lessons back
  \item \emph{Recall:} say \emph{vero?} — \textbf{R3}, durable, twenty to sixty lessons back
  \item \emph{Recall:} say \emph{nero, bianco} — \textbf{R4}, recognition at distance, eighty lessons back or more
\end{itemize}
\end{quote}

\subsection*{You'll want to know: the pair}
\begin{quote}
\textbf{qui} — \textbf{here} · \textbf{lì} — \textbf{there}
\end{quote}

\begin{quote}
\textbf{Il caffè è qui.} — The coffee is here. \textbf{Marco è lì.} — Marco is there. \textbf{Sono qui.} — I'm here.
\end{quote}

Two ways, matching the two ways of pointing:

\noindent
\begin{tabularx}{\linewidth}{@{}>{\raggedright\arraybackslash}X>{\raggedright\arraybackslash}X@{}}
\toprule
\textbf{pointing at a thing} & \textbf{pointing at a place} \\
\midrule
\textbf{questo} — this, near me & \textbf{qui} — here \\
\textbf{quello} — that, away & \textbf{lì} — there \\
\bottomrule
\end{tabularx}

Italian also has \textbf{qua} and \textbf{là}, which mean the same and are used equally often; \emph{qui/lì} are slightly more precise and \emph{qua/là} slightly vaguer, and no Italian will notice which you pick.

\begin{cousinweb}
This is the third lesson in a row built on the same two Latin words.

\textbf{Qui} is \textbf{eccu(m) hīc} — \textquotedblleft{}behold, here.\textquotedblright{} That \emph{hīc} is Latin's near-pointer, the one Italian did \textbf{not} keep for things; it survives only in the place word.

\textbf{Lì} is \textbf{illīc} — \textquotedblleft{}there\textquotedblright{}, and \emph{illīc} is \emph{ille} again, the same far-pointer that gave you \textbf{quello} one lesson ago and \textbf{il, la, lo} on the first page.

So: \emph{ille} became the article, and \emph{quello}, and now \textbf{lì}. One Latin word, three Italian jobs, and this book has now taught all three.

The accent on \textbf{lì} is doing the work it did on \emph{sì} and \emph{né}: it keeps the word apart from \textbf{li}, an unstressed word with another job.
\end{cousinweb}

\subsection*{Your turn}
Say these aloud:

\begin{itemize}
\raggedright
  \item \emph{qui, lì}
  \item \emph{Il caffè è qui.}
  \item \emph{Marco è lì.}
  \item \emph{qua, là}
\end{itemize}

\begin{tcolorbox}[breakable,colback=teal!4,colframe=teal!35!black,title={Before you move on}]
Say \textquotedblleft{}here\textquotedblright{} and \textquotedblleft{}there\textquotedblright{}. (\emph{\textbf{Qui}} — \emph{\textbf{lì}}.) Which Latin word is \emph{lì} built on? (\emph{\textbf{Ille}}, through \emph{illīc}.) Name the three Italian jobs \emph{ille} has done in this book. (\textbf{The article} \emph{\textbf{il/la/lo}}, \textbf{the far demonstrative} \emph{\textbf{quello}}, \textbf{and} \emph{\textbf{lì}}.)
\end{tcolorbox}
